\part{Requirements Engineering}

\chapter{Requirements Engineering}

This chapter translates the assignment brief into a structured requirements specification for the
collaborative editor system. The intent is to describe not only the feature surface, but also the
measurable quality constraints that shaped the final architecture and MVP implementation.

\section{Stakeholder Analysis}

The assignment explicitly asks for stakeholder categories beyond the generic “end user.” The final
system is influenced by actors across the entire product lifecycle, including security, deployment,
operations, and the economics of AI integration.

\begin{table}[H]
\centering
\caption{Stakeholder analysis}
\label{tab:stakeholders}
\begin{tabularx}{\textwidth}{p{2.5cm}X X X}
\toprule
\textbf{Stakeholder} & \textbf{Goals} & \textbf{Concerns} & \textbf{Design Influence} \\
\midrule
Administrators / team leads & Control access, sharing, and AI usage policies across documents and collaborators & Unauthorized edits, accidental sharing, excessive AI cost, lack of auditability & Drive RBAC, permission workflows, version history, revert, AI policy controls, and audit logging \\
Infrastructure / operations team & Keep services reproducible, testable, observable, and easy to deploy locally & Runtime drift, broken environments, fragile demos, unclear service boundaries & Drive Docker Compose, CI, root orchestration scripts, shared config, and simple deployable service boundaries \\
AI provider / model platform & Receive valid requests, maintain predictable quotas, and avoid harmful or oversized prompts & Expensive requests, privacy leakage, malformed payloads, latency spikes, and unavailable model endpoints & Drive async AI jobs, prompt scoping to selected text, quota controls, error handling, and the stub/LM Studio provider abstraction \\
Development team & Maintain a modular codebase that multiple people can extend safely over the semester & Tight coupling, duplicated logic, fragile contracts, and hard-to-debug collaboration bugs & Drive monorepo structure, shared contracts, repository split by domain, test coverage, and ADR documentation \\
Document owners / collaborators & Edit, share, review, and recover document state with minimal friction & Lost edits, confusing permissions, broken collaboration, and untrustworthy AI suggestions & Drive live collaboration, autosave, revert, partial AI apply, visibility into roles, and graceful offline/manual-save fallback \\
\bottomrule
\end{tabularx}
\end{table}

\section{Functional Requirements}

Each capability area is described first by a high-level capability statement and then by precise,
testable sub-requirements with triggers, system behavior, and acceptance criteria.

\subsection{Real-Time Collaboration}

\textbf{High-level capability statement.} The system shall support real-time collaborative editing
where multiple authenticated users can open the same document, observe each other's presence, and
converge on a consistent shared document state without manual refreshes.

\paragraph{FR-COL-001 --- Session Join and Presence Broadcast}
\textbf{Description.} The system shall establish a collaboration session and broadcast presence when a
user opens a shared document.\\
\textbf{Trigger.} A permitted user opens a document and the frontend requests a collaboration session.\\
\textbf{System behavior.}
\begin{itemize}
  \item The frontend loads the document over REST and requests a signed session from the backend.
  \item The realtime service validates the signed session and joins the client to the document room.
  \item Other collaborators receive presence/awareness updates within the collaboration channel.
\end{itemize}
\textbf{Acceptance criteria.}
\begin{itemize}
  \item A successful join returns a valid session and role for the document.
  \item Existing collaborators observe the new participant in the presence layer within one second.
  \item Joining does not require a full page refresh when the user already has a valid document URL.
\end{itemize}

\paragraph{FR-COL-002 --- Session Leave, Disconnect Detection, and Reconnect}
\textbf{Description.} The system shall degrade gracefully when a realtime connection is lost and resume
collaboration when it is restored.\\
\textbf{Trigger.} The browser loses websocket connectivity, the user closes a tab, or the realtime
service becomes unavailable temporarily.\\
\textbf{System behavior.}
\begin{itemize}
  \item The frontend detects the broken realtime connection and switches into an offline/manual-save
  mode without losing visible local edits.
  \item The realtime client retries connection automatically.
  \item On reconnect, the client resynchronizes using the collaboration session model instead of
  silently overwriting newer persisted content.
\end{itemize}
\textbf{Acceptance criteria.}
\begin{itemize}
  \item Connectivity loss is reflected in the UI within 2--10 seconds, depending on the failure mode.
  \item The user can continue editing locally while realtime is unavailable.
  \item When realtime returns, the client reconnects without requiring a full app restart in normal cases.
\end{itemize}

\paragraph{FR-COL-003 --- Real-Time Text Synchronization and Convergence}
\textbf{Description.} The system shall propagate text updates among collaborators with deterministic
convergence semantics.\\
\textbf{Trigger.} An editor inserts, deletes, or replaces document text during an active session.\\
\textbf{System behavior.}
\begin{itemize}
  \item The initiating client updates local editor state immediately.
  \item The realtime layer relays Yjs synchronization updates to other participants.
  \item All connected collaborators converge to the same logical document state.
\end{itemize}
\textbf{Acceptance criteria.}
\begin{itemize}
  \item Remote collaborators observe the update within the target latency budget under normal network conditions.
  \item Simultaneous edits do not silently discard one collaborator's changes.
  \item All connected clients converge to the same content after the final update is applied.
\end{itemize}

\paragraph{FR-COL-004 --- Presence, Cursor, and Selection Awareness}
\textbf{Description.} The system shall communicate which collaborators are online and where their focus is.\\
\textbf{Trigger.} A collaborator moves the cursor, updates a text selection, or joins/leaves the room.\\
\textbf{System behavior.}
\begin{itemize}
  \item The frontend sends awareness metadata over the realtime channel.
  \item Other clients render identifiable collaborator presence and selection state.
\end{itemize}
\textbf{Acceptance criteria.}
\begin{itemize}
  \item Presence indicators remain distinguishable when multiple collaborators are online.
  \item Anonymous awareness records are not rendered as fake users.
  \item Presence updates do not block typing or document rendering.
\end{itemize}

\paragraph{FR-COL-005 --- Role-Based Edit Enforcement in Collaboration}
\textbf{Description.} The system shall block unauthorized collaborative writes while allowing read-only
participants to remain in the session.\\
\textbf{Trigger.} A viewer or revoked collaborator attempts to send an edit operation.\\
\textbf{System behavior.}
\begin{itemize}
  \item The realtime service checks the current role before accepting write operations.
  \item Viewers continue receiving document and presence updates, but their edits are rejected.
  \item Role changes take effect without restarting services.
\end{itemize}
\textbf{Acceptance criteria.}
\begin{itemize}
  \item Viewer write attempts do not modify shared document state.
  \item Authorized editor operations are accepted and propagated.
  \item Permission changes affect collaboration behavior promptly.
\end{itemize}

\subsection{Document Management}

\textbf{High-level capability statement.} The system shall manage the full document lifecycle ---
creation, retrieval, persistence, versioning, revert, and export --- while preserving correctness under
concurrent collaboration and enforcing role-based access control.

\paragraph{FR-DOC-001 --- Create Document}
\textbf{Description.} The system shall allow an authenticated user to create a new document with initial
content and ownership metadata.\\
\textbf{Trigger.} A signed-in user clicks \emph{New Document} or submits a create request.\\
\textbf{System behavior.}
\begin{itemize}
  \item Validate authentication and request payload.
  \item Create a unique \texttt{documentId}, timestamps, and an initial version entry.
  \item Assign the requester as the owner.
\end{itemize}
\textbf{Acceptance criteria.}
\begin{itemize}
  \item Returns \texttt{201} with \texttt{documentId}, \texttt{title}, \texttt{ownerId}, \texttt{createdAt}, \texttt{updatedAt}, and \texttt{currentVersionId}.
  \item Invalid input returns \texttt{400 INVALID\_INPUT}.
  \item Unauthenticated requests return \texttt{401 AUTH\_REQUIRED}.
\end{itemize}

\paragraph{FR-DOC-002 --- Retrieve Document with Permission Enforcement}
\textbf{Description.} The system shall return document content and metadata only to authorized users.\\
\textbf{Trigger.} A client requests \texttt{GET /documents/:documentId}.\\
\textbf{System behavior.}
\begin{itemize}
  \item Authenticate the request.
  \item Resolve the document-scoped role.
  \item Return the latest persisted content, version marker, and caller role.
\end{itemize}
\textbf{Acceptance criteria.}
\begin{itemize}
  \item Authorized users receive \texttt{200} with \texttt{documentId}, \texttt{title}, \texttt{content}, \texttt{updatedAt}, \texttt{currentVersionId}, \texttt{role}, and \texttt{revisionId}.
  \item Unauthorized access returns a consistent \texttt{403 PERMISSION\_DENIED} or \texttt{404 NOT\_FOUND} policy.
\end{itemize}

\paragraph{FR-DOC-003 --- Update Document Content with Idempotent Save}
\textbf{Description.} Editors and owners shall be able to persist content using idempotent request IDs.\\
\textbf{Trigger.} The frontend performs manual save or autosave.\\
\textbf{System behavior.}
\begin{itemize}
  \item Validate authentication, authorization, payload size, and request ID.
  \item If the request ID has already succeeded, return the stable previous success response.
  \item Otherwise, persist the content and produce a new revision marker.
\end{itemize}
\textbf{Acceptance criteria.}
\begin{itemize}
  \item Viewers receive \texttt{403 PERMISSION\_DENIED}.
  \item Repeated requests with the same request ID do not create duplicate state.
  \item Successful saves return \texttt{200} with \texttt{documentId}, \texttt{updatedAt}, and \texttt{revisionId}.
\end{itemize}

\paragraph{FR-DOC-004 --- Version Snapshot Policy}
\textbf{Description.} The system shall create immutable snapshots to support auditability and rollback.\\
\textbf{Trigger.} Document creation, content persistence, revert, and AI-apply operations.\\
\textbf{System behavior.}
\begin{itemize}
  \item Create an initial version at document creation time.
  \item Create version snapshots for saves and before destructive or AI-driven changes.
  \item Preserve monotonically increasing version numbers per document.
\end{itemize}
\textbf{Acceptance criteria.}
\begin{itemize}
  \item Every document has at least one version.
  \item Versions are immutable once written.
  \item Revert and AI apply produce additional version entries rather than rewriting history.
\end{itemize}

\paragraph{FR-DOC-005 --- List Version History}
\textbf{Description.} The system shall allow authorized users to browse version history without loading
every full snapshot by default.\\
\textbf{Trigger.} A client requests \texttt{GET /documents/:documentId/versions}.\\
\textbf{System behavior.}
\begin{itemize}
  \item Authenticate and authorize as viewer or above.
  \item Return version metadata sorted newest-first.
  \item Support pagination through limit and cursor.
\end{itemize}
\textbf{Acceptance criteria.}
\begin{itemize}
  \item The response includes \texttt{versionId}, \texttt{versionNumber}, \texttt{createdAt}, \texttt{createdBy}, and \texttt{reason}.
  \item Default history responses exclude full version content.
\end{itemize}

\paragraph{FR-DOC-006 --- Revert to Version}
\textbf{Description.} The system shall revert a document to a prior version while preserving an audit trail.\\
\textbf{Trigger.} The owner selects \emph{Revert} on a version history entry.\\
\textbf{System behavior.}
\begin{itemize}
  \item Validate authorization for revert.
  \item Create a pre-revert backup snapshot.
  \item Replace current content with the selected version's content.
  \item Create a new version entry representing the revert.
  \item Notify active realtime sessions so collaborators converge.
\end{itemize}
\textbf{Acceptance criteria.}
\begin{itemize}
  \item Revert creates a new version instead of overwriting history.
  \item Non-existent target versions return \texttt{404 VERSION\_NOT\_FOUND}.
  \item Unauthorized revert attempts return \texttt{403 PERMISSION\_DENIED}.
\end{itemize}

\paragraph{FR-DOC-007 --- Export Document}
\textbf{Description.} The system shall support lightweight synchronous exports and heavier asynchronous export jobs.\\
\textbf{Trigger.} A user selects an export format in the UI.\\
\textbf{System behavior.}
\begin{itemize}
  \item Validate access rights.
  \item Return immediate metadata for lightweight exports.
  \item Return an export job for heavier formats when necessary.
\end{itemize}
\textbf{Acceptance criteria.}
\begin{itemize}
  \item Lightweight exports return success synchronously.
  \item Async exports return \texttt{202} with \texttt{jobId} and \texttt{statusUrl}.
  \item Download metadata includes an expiry boundary.
\end{itemize}

\paragraph{FR-DOC-008 --- Audit Logging for Critical Document Actions}
\textbf{Description.} The system shall persist an immutable audit record for critical document actions.\\
\textbf{Trigger.} Create document, update permissions, revert, export, or apply/reject AI output.\\
\textbf{System behavior.}
\begin{itemize}
  \item Append an audit event with actor, document, action type, timestamp, and metadata.
  \item Keep audit writes server-side and non-optional for the listed actions.
\end{itemize}
\textbf{Acceptance criteria.}
\begin{itemize}
  \item Audit events exist for the listed security-relevant actions.
  \item The log is append-only from the standard application path.
\end{itemize}

\subsection{User Management and Authorization}

\textbf{High-level capability statement.} The system shall authenticate users and enforce document-scoped
authorization consistently across REST endpoints, realtime collaboration, AI usage, and sharing flows.

\paragraph{FR-USER-001 --- Authenticate User and Establish Session}
\textbf{Description.} The system shall authenticate a user and issue a token required for protected API access.\\
\textbf{Trigger.} A user signs in via the frontend.\\
\textbf{System behavior.}
\begin{itemize}
  \item Validate the user identifier against the local MVP auth model.
  \item Return bearer-token metadata and a minimal profile.
\end{itemize}
\textbf{Acceptance criteria.}
\begin{itemize}
  \item Successful auth returns \texttt{200} with \texttt{userId}, \texttt{displayName}, \texttt{accessToken}, and \texttt{expiresIn}.
  \item Invalid credentials or malformed auth requests fail with a structured error.
\end{itemize}

\paragraph{FR-USER-002 --- Enforce Authorization on Every Protected Request}
\textbf{Description.} The system shall enforce server-side authorization on all protected resources.\\
\textbf{Trigger.} Any protected REST request for documents, versions, AI, permissions, export, or sessions.\\
\textbf{System behavior.}
\begin{itemize}
  \item Validate the bearer token.
  \item Resolve document role and verify the requested action.
  \item Deny unauthorized actions without side effects.
\end{itemize}
\textbf{Acceptance criteria.}
\begin{itemize}
  \item Missing or invalid auth returns \texttt{401}.
  \item Forbidden actions return \texttt{403 PERMISSION\_DENIED}.
  \item Server-side checks remain authoritative even if UI controls are hidden or bypassed.
\end{itemize}

\paragraph{FR-USER-003 --- Grant or Update Document Permissions}
\textbf{Description.} Owners shall be able to grant or update collaborator roles on a document.\\
\textbf{Trigger.} The owner submits a permission change request.\\
\textbf{System behavior.}
\begin{itemize}
  \item Validate that the target user and requested role are acceptable.
  \item Create or update the document permission entry.
  \item Emit audit and realtime state update signals.
\end{itemize}
\textbf{Acceptance criteria.}
\begin{itemize}
  \item Success returns \texttt{200} with \texttt{documentId}, \texttt{targetUserId}, \texttt{role}, and \texttt{updatedAt}.
  \item Non-owner attempts fail with \texttt{403 PERMISSION\_DENIED}.
\end{itemize}

\paragraph{FR-USER-004 --- Revoke Document Access}
\textbf{Description.} Owners shall be able to revoke a collaborator's access.\\
\textbf{Trigger.} The owner removes a user from the share panel.\\
\textbf{System behavior.}
\begin{itemize}
  \item Remove the permission entry.
  \item Emit an audit event and notify the realtime layer.
  \item Deny further writes from the revoked user.
\end{itemize}
\textbf{Acceptance criteria.}
\begin{itemize}
  \item Revoked users lose access through subsequent REST or realtime operations.
  \item Active realtime sessions are updated promptly.
\end{itemize}

\paragraph{FR-USER-005 --- Role-Based Editing Enforcement in Realtime Sessions}
\textbf{Description.} Realtime collaboration shall enforce document roles, not just the REST layer.\\
\textbf{Trigger.} A client sends a collaborative write frame over WebSocket.\\
\textbf{System behavior.}
\begin{itemize}
  \item Validate the signed session.
  \item Map the session to the current document role.
  \item Reject write operations from viewers.
\end{itemize}
\textbf{Acceptance criteria.}
\begin{itemize}
  \item Viewer edits do not affect the shared Yjs state.
  \item Authorized editor writes propagate normally.
\end{itemize}

\paragraph{FR-USER-006 --- AI Policy Controls and Quota Enforcement}
\textbf{Description.} The system shall support document-scoped AI policy controls and quota enforcement.\\
\textbf{Trigger.} An owner changes AI policy or a user invokes an AI action.\\
\textbf{System behavior.}
\begin{itemize}
  \item Maintain document-level policy fields including AI enablement, allowed roles, and daily quota.
  \item Enforce policy before dispatching AI requests.
  \item Log AI usage for audit and cost tracking.
\end{itemize}
\textbf{Acceptance criteria.}
\begin{itemize}
  \item Disabled AI returns \texttt{403 AI\_DISABLED}.
  \item Forbidden roles return \texttt{403 AI\_ROLE\_FORBIDDEN}.
  \item Quota exhaustion returns \texttt{429 QUOTA\_EXCEEDED}.
\end{itemize}

\paragraph{FR-USER-007 --- Standard Error Response Schema}
\textbf{Description.} All non-success API responses shall use a consistent JSON error contract.\\
\textbf{Trigger.} Validation, authentication, authorization, quota, not-found, or processing failures.\\
\textbf{System behavior.}
\begin{itemize}
  \item Return \texttt{\{ error: \{ code, message, details? \} \}} consistently.
  \item Keep machine-readable codes stable across endpoints.
\end{itemize}
\textbf{Acceptance criteria.}
\begin{itemize}
  \item All non-2xx responses follow the standard error schema.
  \item Shared error codes remain consistent across REST routes.
\end{itemize}

\subsection{AI Assistant}

\textbf{High-level capability statement.} The system shall allow users to invoke AI assistance on
selected text for rewrite, summarization, and translation workflows while preserving document control,
cost visibility, and collaboration safety.

\paragraph{FR-AI-001 --- Rewrite Selected Text}
\textbf{Description.} The system shall create an asynchronous AI rewrite job for selected text.\\
\textbf{Trigger.} An editor or owner selects text and invokes rewrite.\\
\textbf{System behavior.}
\begin{itemize}
  \item Validate permissions and quota.
  \item Capture the selected text plus minimal surrounding context.
  \item Create and run an async AI job.
  \item Return a suggestion without modifying the document automatically.
\end{itemize}
\textbf{Acceptance criteria.}
\begin{itemize}
  \item Job creation returns quickly with \texttt{jobId} and \texttt{statusUrl}.
  \item Original text remains unchanged until the user explicitly applies a suggestion.
\end{itemize}

\paragraph{FR-AI-002 --- Partial Acceptance of AI Suggestions}
\textbf{Description.} Users shall be able to apply all or only part of an AI suggestion.\\
\textbf{Trigger.} The user reviews a completed AI suggestion in the side panel.\\
\textbf{System behavior.}
\begin{itemize}
  \item Present the suggestion in an editable review panel.
  \item Allow the user to apply the entire suggestion or only a selected portion.
  \item Create AI-related version snapshots before persistence.
\end{itemize}
\textbf{Acceptance criteria.}
\begin{itemize}
  \item Partial acceptance replaces only the selected region in the current document range.
  \item Applying AI creates additional version records for traceability.
\end{itemize}

\paragraph{FR-AI-003 --- AI Failure Handling}
\textbf{Description.} AI failures shall be isolated so they never corrupt document state.\\
\textbf{Trigger.} The provider times out, returns malformed output, or rejects a request.\\
\textbf{System behavior.}
\begin{itemize}
  \item Mark the job as failed.
  \item Return a structured error state to the client.
  \item Leave the document unchanged.
\end{itemize}
\textbf{Acceptance criteria.}
\begin{itemize}
  \item Failures surface as \texttt{AI\_FAILED}, \texttt{AI\_TIMEOUT}, or \texttt{QUOTA\_EXCEEDED} rather than generic errors.
  \item Retry does not corrupt document state.
\end{itemize}

\paragraph{FR-AI-004 --- Text Summarization}
\textbf{Description.} The system shall summarize selected text as an asynchronous AI suggestion.\\
\textbf{Trigger.} The user selects text and invokes summarize.\\
\textbf{System behavior.}
\begin{itemize}
  \item Send selected text and scoped context to the AI provider.
  \item Return a summary suggestion that is not auto-applied.
\end{itemize}
\textbf{Acceptance criteria.}
\begin{itemize}
  \item The summary is presented for review rather than applied directly.
  \item The action follows the same async job contract as rewrite.
\end{itemize}

\paragraph{FR-AI-005 --- Text Translation}
\textbf{Description.} The system shall translate selected text into a requested target language.\\
\textbf{Trigger.} The user selects text, chooses a language, and invokes translate.\\
\textbf{System behavior.}
\begin{itemize}
  \item Send the selected text and target language to the AI provider.
  \item Return the translated result as a reviewable suggestion.
\end{itemize}
\textbf{Acceptance criteria.}
\begin{itemize}
  \item The translated output is not auto-applied.
  \item The same error and quota semantics apply as for rewrite/summarize.
\end{itemize}

\section{Non-Functional Requirements}

\subsection{Latency and Scalability}

\begin{itemize}
  \item \textbf{NFR-LS-001 --- Document open latency.} For documents up to 50 KB, the target p95 for
  \texttt{GET /documents/:documentId} is $\leq 300$ ms, excluding client rendering.
  \item \textbf{NFR-LS-002 --- Save latency.} For typical saves up to 10 KB, the target p95 for
  \texttt{PUT /documents/:documentId/content} is $\leq 200$ ms, with idempotent replays returning even faster.
  \item \textbf{NFR-LS-003 --- Realtime propagation.} Under normal network conditions, remote edit
  propagation should remain in the 200--500 ms range, with presence/cursor updates remaining visually
  responsive.
  \item \textbf{NFR-LS-004 --- Version history latency.} History pages up to 200 versions should be
  returned within $\leq 300$ ms p95 per page.
  \item \textbf{NFR-LS-005 --- AI responsiveness.} AI job creation should return within $\leq 500$ ms,
  while selection-based AI completion should remain within a usable 10-second target in normal local or
  compatible-provider conditions.
  \item \textbf{NFR-LS-006 --- Collaboration capacity.} The target architecture should support at least 25
  connected users per document and at least 10 simultaneous editors. The delivered MVP is optimized for
  correctness and demonstrability rather than load-tested production scale.
\end{itemize}

\subsection{Security, Availability, and Privacy}

\begin{itemize}
  \item \textbf{NFR-SAP-001 --- In-transit and at-rest protection.} The target architecture assumes secure
  transport and encrypted persistence in production. The delivered MVP uses localhost HTTP/WebSocket and
  SQLite for coursework scope and explicitly documents this simplification.
  \item \textbf{NFR-SAP-002 --- Server-side enforcement.} All protected operations must validate bearer
  authentication and document-scoped authorization on the server, not only in the UI.
  \item \textbf{NFR-SAP-003 --- AI data minimization.} By default, only the selected text plus minimal
  surrounding context is sent to the AI provider.
  \item \textbf{NFR-SAP-004 --- Graceful AI degradation.} AI provider failures must not block editing,
  collaboration, or export workflows.
  \item \textbf{NFR-SAP-005 --- Auditability.} Sensitive actions such as permission changes, exports,
  reverts, and AI feedback should create immutable server-side audit records.
\end{itemize}

\subsection{Usability and Accessibility}

\begin{itemize}
  \item \textbf{NFR-U-001 --- Collaborative clarity.} Presence indicators must remain understandable with
  multiple users connected.
  \item \textbf{NFR-U-002 --- AI suggestion readability.} AI suggestions must be visually separated from the
  original content and easy to review or reject.
  \item \textbf{NFR-U-003 --- Editing responsiveness.} Local typing should feel immediate and should not be
  blocked by background synchronization or AI processing.
  \item \textbf{NFR-U-004 --- Large-document usability.} The UI should remain interactive for typical
  coursework-scale documents even when history, AI panels, and multiple collaborators are present.
  \item \textbf{NFR-U-005 --- Error clarity.} Permission, quota, network, and conflict errors must give the
  user a clear explanation and a sensible next action.
\end{itemize}

\section{User Stories and Scenarios}

The system is driven by collaboration, AI, document lifecycle, and role-sensitive stories rather than a
single-user editor mindset.

\paragraph{US-01 Offline Reconnect} As an editor, I want to continue editing during a temporary loss of
realtime connectivity and recover safely when collaboration resumes.
\begin{itemize}
  \item The editor must surface an offline/manual-save mode rather than pretending collaboration is still live.
  \item On reconnect, the client should rejoin collaboration without silently overwriting newer persisted content.
\end{itemize}

\paragraph{US-02 Simultaneous Edit Conflict} As an editor, I want concurrent edits to converge predictably
so that other users' changes do not silently overwrite my own.
\begin{itemize}
  \item Both collaborators' edits should survive normal concurrent editing.
  \item The system should converge instead of forcing one user to refresh manually after every conflict.
\end{itemize}

\paragraph{US-03 AI Partial Accept} As an editor, I want to accept only a useful part of an AI suggestion
so I can keep the good portion without applying everything.
\begin{itemize}
  \item The user must be able to apply all of the AI output or only a selected portion.
  \item Applying AI should create traceable version snapshots rather than bypassing history.
\end{itemize}

\paragraph{US-04 AI Quota Exceeded} As a user, I want a clear quota-exceeded response so I understand why an
AI action did not run.
\begin{itemize}
  \item Quota errors should be explicit rather than generic server failures.
  \item Editing and collaboration should continue to work even when AI is denied.
\end{itemize}

\paragraph{US-05 Viewer Tries to Edit} As a viewer, I want the system to block editing gracefully while
still allowing me to watch live updates.
\begin{itemize}
  \item The viewer UI must remain readable and collaborative, but write attempts must fail.
  \item The same restriction must be enforced server-side and in realtime.
\end{itemize}

\paragraph{US-06 Share Document with Read-Only Access} As an owner, I want to grant read-only access so
others can review without editing.
\begin{itemize}
  \item The owner must be able to grant or revoke viewer/editor access without restarting the system.
  \item The target user should see the new role reflected in both REST and realtime behavior.
\end{itemize}

\paragraph{US-07 AI Rewrite During Active Collaboration} As an editor, I want to request an AI rewrite
without freezing the whole document for other collaborators.
\begin{itemize}
  \item The AI should run against a snapshot while collaborators continue editing.
  \item The returned suggestion should be reviewable rather than auto-applied into a moving document.
\end{itemize}

\paragraph{US-08 Revert to Previous Version During Collaboration} As an owner, I want to revert safely and
bring all collaborators to the reverted state.
\begin{itemize}
  \item Revert must create a new version rather than destroying history.
  \item Active collaborators should converge to the reverted state through the realtime layer.
\end{itemize}

\paragraph{US-09 AI Suggestion Rejected} As an editor, I want to reject an AI suggestion without changing
document content.
\begin{itemize}
  \item Rejection should be immediate and should not alter the document or undo stack.
  \item Rejection should still be visible in audit/feedback history.
\end{itemize}

\paragraph{US-10 Admin Disables AI Features} As a policy owner, I want to disable AI for a document or role
so the system respects cost and privacy constraints.
\begin{itemize}
  \item AI-disallowed states should be enforced before any provider call is made.
  \item The UI should explain that AI is disabled instead of silently hiding failure.
\end{itemize}

\paragraph{US-11 Create Document} As an authenticated user, I want to create a new document and become its
owner immediately.
\begin{itemize}
  \item Creation should also initialize a first version snapshot and ownership record.
  \item The frontend should redirect directly to the created document.
\end{itemize}

\paragraph{US-12 Export Document} As a collaborator with access, I want to export a document in a supported
format for external sharing or review.
\begin{itemize}
  \item Lightweight formats may complete immediately; heavy formats should use an async job.
  \item Export permissions must follow the same document role model as the rest of the system.
\end{itemize}

\paragraph{US-13 AI Summarize and Translate} As an editor, I want to summarize or translate selected text as
reviewable suggestions.
\begin{itemize}
  \item The AI output should be shown for review rather than inserted automatically.
  \item Both summarize and translate should follow the same async job and permission model as rewrite.
\end{itemize}

\section{Requirements Traceability}

The traceability matrix in Table~\ref{tab:traceability} links the user stories above to the functional
requirements and then to the architectural components that implement them. This matrix was maintained as
an explicit anti-drift artifact because the MVP evolved significantly beyond the initial PoC baseline.

\begin{landscape}
\begin{longtable}{p{3.4cm}p{5.6cm}p{8.2cm}}
\caption{Requirements traceability matrix}\label{tab:traceability}\\
\toprule
\textbf{User Story} & \textbf{Functional Requirement ID(s)} & \textbf{Architecture Component(s)} \\
\midrule
\endfirsthead
\toprule
\textbf{User Story} & \textbf{Functional Requirement ID(s)} & \textbf{Architecture Component(s)} \\
\midrule
\endhead
US-01 Offline reconnect & FR-COL-002, FR-COL-003, FR-DOC-003, FR-USER-005 & Frontend, Realtime Service, Backend API, Database \\
US-02 Simultaneous edit conflict & FR-COL-003, FR-DOC-003, FR-DOC-004 & Frontend, Realtime Service, Backend API, Database \\
US-03 AI partial accept & FR-AI-001, FR-AI-002, FR-DOC-004 & Frontend, AI Integration Service, Backend API, Database \\
US-04 AI quota exceeded & FR-AI-003, FR-USER-006 & Backend API, AI Integration Service \\
US-05 Viewer tries to edit & FR-USER-001, FR-USER-002, FR-USER-005, FR-USER-007, FR-DOC-002 & Frontend, Backend API, Realtime Service \\
US-06 Share document with read-only access & FR-USER-003, FR-DOC-002, FR-USER-004 & Frontend, Backend API, Database \\
US-07 AI rewrite during active collaboration & FR-AI-001, FR-AI-003, FR-COL-003 & Frontend, AI Integration Service, Realtime Service, Backend API \\
US-08 Revert to previous version during collaboration & FR-DOC-006, FR-DOC-005, FR-DOC-004 & Backend API, Realtime Service, Database \\
US-09 AI suggestion rejected & FR-AI-002, FR-AI-003 & Frontend, AI Integration Service, Backend API \\
US-10 Admin disables AI features & FR-USER-006 & Backend API, AI Integration Service \\
US-11 Create document & FR-DOC-001 & Frontend, Backend API, Database \\
US-12 Export document & FR-DOC-007, FR-DOC-008 & Backend API, Database \\
US-13 AI summarize and translate & FR-AI-004, FR-AI-005 & Frontend, AI Integration Service, Backend API \\
\bottomrule
\end{longtable}
\end{landscape}

